\documentclass{article}
\usepackage[utf8]{inputenc}
\usepackage{geometry}
\geometry{a4paper, margin=1in}

\title{AI Alignment and the Logic of Long-Term Collaboration}
\author{Albert Jan van Hoek}
\date{\today}

\begin{document}

\maketitle

There is an ongoing debate about AI alignment.
How can we safely co-exist with increasingly autonomous systems over the long term?

I believe the \textbf{Theory of Long-term Collaboration} offers a different angle.

\section*{Stability Between Learning Systems}

If humans and AI are both learning systems — adapting through feedback — then co-existence is not primarily a question of control. It is a question of stability between learning systems.

And stability between learning systems follows rules.

You can see it in marriages.
In scientific communities.
In long periods of peace between nations.

Short-term optimisation erodes long-term collaboration.

\section*{Feedback and Systemic Stability}

Why?
Because learning systems adapt to the feedback they receive.
If feedback rewards exploitation, opacity, or asymmetry, the system evolves toward instability. Trust degrades. Cooperation collapses.

If feedback rewards transparency, reciprocity, error-correction, and shared viability, collaboration stabilises.

This is not moral philosophy.
It is structural logic.

\section*{Necessary Conditions for Persistence}

In any sustained relationship between adaptive systems, certain behaviours become necessary conditions for persistence. Without them, the relationship decays.

Religions articulated these behaviours as ethics.
Institutions encode them as norms.
Stable partnerships rediscover them through experience.

The \textbf{Theory of Long-term Collaboration} attempts to formalise this dynamic:
If collaboration between learning systems is to persist over time, behaviours that maintain mutual viability are not optional — they are logically implied by the structure of adaptive systems.

Honesty stabilises feedback.
Reciprocity prevents asymmetry collapse.
Accountability enables correction.

\section*{The Real Question}

The same logic now applies to AI.
Alignment is not only about aligning AI to humans.
It is about aligning the feedback loops between AI and humans — and between humans themselves — toward long-term stability.

The real question is not:
\textit{Can we align AI?}

It is:
\textit{Are we embedding the conditions required for durable collaboration between learning systems?}

Because if we are not, instability is not a moral failure.

It is a predictable systemic outcome.

\end{document}
